% Generated draft: requires mathematical and source review.
\documentclass[11pt]{article}
\usepackage[margin=25mm]{geometry}
\usepackage{fontspec}
\setmainfont{Latin Modern Roman}
\usepackage{amsmath,amssymb,mathtools,unicode-math}
\setmathfont{Latin Modern Math}
\usepackage{enumitem,needspace,xurl,hyperref}
\setlength{\parindent}{0pt}
\setlength{\parskip}{5pt}
\setlength{\emergencystretch}{3em}
\newcommand{\N}{\mathbb{N}}
\newcommand{\Z}{\mathbb{Z}}
\newcommand{\Q}{\mathbb{Q}}
\newcommand{\R}{\mathbb{R}}
\newcommand{\C}{\mathbb{C}}
\newcommand{\F}{\mathbb{F}}
\newcommand{\A}{\mathbb{A}}
\newcommand{\PP}{\mathbb{P}}
\newcommand{\E}{\mathbb{E}}
\newcommand{\eps}{\varepsilon}
\newcommand{\dd}{\,\mathrm{d}}
\begin{document}
\section*{UnsolvedMath ID: 1100402}
\textbf{Source problem code:} catalogue entry
\subsection{Definitions and mathematical statement}
Let $S^1=\{z\in\mathbb C:|z|=1\}$ and let $\operatorname{Homeo}^{+}(S^1)$ be the group of orientation-preserving homeomorphisms of the circle, with the compact-open topology. For an integer $k\geq1$, let
\[
\mathcal C_k(S^1)=\{A\subset S^1:|A|=k\}
\]
be the space of unordered configurations of $k$ distinct points. The induced action of a subgroup $G$ is $g\cdot A=\{g(a):a\in A\}$. It is transitive if for every $A,B\in\mathcal C_k(S^1)$ there exists $g\in G$ such that $g(A)=B$.

\textbf{Primary question.} Does there exist a subgroup $G$ such that
\[
G\subsetneq\operatorname{Homeo}^{+}(S^1),\qquad
G\text{ is closed},\qquad
G\text{ acts transitively on }\mathcal C_4(S^1)?
\]
\textbf{Further questions.} For which positive integers $k\ne3$ does such a proper closed subgroup acting transitively on $\mathcal C_k(S^1)$ exist? What relation, if any, do these questions have to earthquakes? The source does not specify a proposition or a particular earthquake construction for this last request. Its unordered-tuples wording is read as configurations of distinct points; allowing repetitions would introduce different multiplicity strata.

The supplied research summary reports a negative answer for $k\geq4$ in the category of closed subgroups containing a nonconstant continuous path, attributed to Giblin--Markovic. It does not report a complete resolution without this additional hypothesis. The stated continuous-transitivity results and the claimed implications from ordinary unordered transitivity must not be identified without their additional hypotheses.
\subsection{Short English statement}
Ask whether a proper closed group of orientation-preserving circle homeomorphisms can move any four-point subset to any other, and ask the analogous questions for other $k\ne3$. The supplied progress addresses groups with a nonconstant continuous path and leaves the unrestricted four-point question unresolved.
\subsection{Sources}
\begin{itemize}
\item M. Bestvina, \emph{Questions in Geometric Group Theory} (2004), Question 4.2, PDF page 12: \url{https://www.math.utah.edu/~bestvina/eprints/questions-updated.pdf}; manual transcription, accessed 2026-07-29.
\item Giblin--Markovic, \emph{Geometry \& Topology} (2006), cited in the research summary for classification and continuous-transitivity results; no title, volume, pages, or theorem number is supplied. The summary also names Cameron in an orbit-counting argument, without a bibliographic reference, and mentions de la Harpe in connection with the separate Question 4.1.
\item Local registry \texttt{data/unsolvedmath/problems.json}, permanent ID $1100402$, source code catalogue entry, \emph{Questions in Geometric Group Theory --- Q 4.2}. The record labels this partially solved and PARTIAL-PROGRESS; its source background requires current-status and rights review.
\end{itemize}
\end{document}
